\begin{frame}{Preparing a ML model}{The overfitting problem\dots}
    \begin{center}
        \def\overfitcurve{(0.30,1.10) (0.37,1.97) (0.45,1.90) (0.52,1.54) (0.59,1.25) (0.67,1.15) (0.74,1.25) (0.81,1.48) (0.89,1.77) (0.96,2.05) (1.03,2.25) (1.11,2.37) (1.18,2.38) (1.25,2.31) (1.33,2.18) (1.40,2.02) (1.47,1.86) (1.55,1.74) (1.62,1.67) (1.69,1.65) (1.77,1.69) (1.84,1.77) (1.92,1.86) (1.99,1.96) (2.06,2.04) (2.14,2.09) (2.21,2.09) (2.28,2.05) (2.36,1.98) (2.43,1.88) (2.50,1.78) (2.58,1.69) (2.65,1.62) (2.72,1.57) (2.80,1.52) (2.87,1.44) (2.94,1.33) (3.02,1.14) (3.09,0.92) (3.16,0.77) (3.20,0.80)}
\def\goodcurve{(0.30,1.05) (0.50,1.31) (0.70,1.54) (0.90,1.72) (1.10,1.87) (1.30,1.98) (1.50,2.04) (1.70,2.06) (1.90,2.05) (2.10,1.99) (2.30,1.88) (2.50,1.74) (2.70,1.55) (2.90,1.31) (3.10,1.03) (3.20,0.88)}
\def\underfitcurve{(0.30,1.69) (3.20,1.54)}
\def\trainpts{0.30/1.10, 0.62/1.19, 0.94/1.99, 1.27/2.29, 1.59/1.70,
    1.91/1.86, 2.23/2.08, 2.56/1.71, 2.88/1.43, 3.20/0.80}
\def\testpts{0.36/0.99, 0.98/1.68, 1.19/2.13, 2.10/2.21,
    2.52/1.85, 2.65/1.47, 2.80/1.53, 2.86/1.63}

\begin{tikzpicture}[
        train/.style={fill=CEAred, draw=none, circle, inner sep=1.6pt},
        test/.style={fill=CEAblue, draw=none, circle, inner sep=1.6pt},
        ax/.style={thick, CEAblack, ->}
    ]
    \foreach \xs/\firsttitle/\subtitle/\title/\curve/\col/\msetrain/\msetest/\name/\effect/\adjective in {
            0/Model 1/Overfitting/High Variance/\overfitcurve/CEAblack/0.0/0.153/overfitting/memorised the training data/unreliable,
            4.7/Model 2/Generalisation/Optimal/\goodcurve/CEAblack/0.044/0.030/generalisation/adjusted to the training data/great,
            9.4/Model 3/Underfitting/High Bias/\underfitcurve/CEAblack/0.200/0.148/underfitting/not learnt the training data/useless}
        {
            \begin{scope}[xshift=\xs cm]
                \draw[ax] (-0.2,0) -- (3.6,0) node[below] {$x$};
                \draw[ax] (0,-0.2) -- (0,3.0) node[left] {$y$};
                \node<1-3>[above] at (1.7,2.9) {\textbf{\firsttitle}};
                \node<4->[above] at (1.7,2.9) {\textcolor{CEAred}{\textbf{\title}}};
                \node<4->[above] at (1.7,3.4) {\footnotesize\textcolor{CEAred}{\textbf{\subtitle}}};
                \draw[\col, thick, smooth] plot coordinates {\curve};
                \foreach \px/\py in \trainpts \node[train] at (\px,\py) {};
                \only<3->{%
                \foreach \px/\py in \testpts \node[test] at (\px,\py) {};
                \node<2->[below, CEAred] at (1.7,-0.4) {$\operatorname{MSE}_{\text{train}}=\msetrain$};
                \node<3->[below, CEAblue] at (1.7,-0.9) {$\operatorname{MSE}_{\text{test}}=\msetest$};
                }

                \node<5->[below, CEAblack, align=center, text width=15em] at (1.7, -1.5) {\footnotesize The \textbf{\name} regime has \emph{\effect}. It is \textbf{\adjective} for new data.};
            \end{scope}
        }
\end{tikzpicture}

    \end{center}

    \vfill

    \begin{onlyenv}<1>
        \includesvg[height=12pt]{01_slides/img/brain.svg}~What would you consider as a \textbf{good model}? Think in terms of something like the \textbf{Mean Squared Error} (MSE):
        \begin{equation*}
            \operatorname{MSE}
            =
            \mathbb{E}\left[ \left( y - f(x) \right)^2 \right]
            =
            \mathbb{E}\left[ \operatorname{dist}\left( y - f(x) \right)^2 \right]
        \end{equation*}
    \end{onlyenv}

    \begin{onlyenv}<4>
        \begin{center}
            The model that seemed to work best, in reality \textbf{does not generalise}!
        \end{center}
    \end{onlyenv}

    \vfill
\end{frame}

\begin{frame}{Preparing a ML model}{How to avoid over- and under-fitting\dots}
    Suppose you have a \textbf{data distribution} $X = \left\{ x_i \right\}_{i = 1}^N$.
    Are there ways to avoid under- and overfitting?

    \begin{center}
        \begin{tikzpicture}
    \draw<1>[thick, draw=CEAred, fill=CEAred!5] (0, 0) rectangle (10cm, 1cm) node[midway] {DATA};
    \draw<2>[thick, draw=CEAred, fill=CEAred!5] (0, 0) rectangle (7cm, 1cm) node[midway] {TRAIN DATA};
    \draw<2>[thick, draw=CEAblue, fill=CEAblue!5] (7cm, 0) rectangle (10cm, 1cm) node[midway] {TEST DATA};
    \draw<3->[thick, draw=CEAred, fill=CEAred!5] (0, 0) rectangle (4cm, 1cm) node[midway] {TRAIN DATA};
    \draw<3->[thick, draw=CEAgreen, fill=CEAgreen!5] (4cm, 0) rectangle (7cm, 1cm) node[midway] {VALIDATION DATA};
    \draw<3->[thick, draw=CEAblue, fill=CEAblue!5] (7cm, 0) rectangle (10cm, 1cm) node[midway] {TEST DATA};
\end{tikzpicture}

    \end{center}

    \begin{itemize}[<+->]
        \item \textcolor{CEAred}{\textbf{Test Set:}} Split \textbf{randomly} your data to keep some of them for \textbf{testing}
        \item \textcolor{CEAred}{\textbf{Validation Set:}} Retain some more data to \textbf{validate} your choice of the ML model
        \item \textcolor{CEAred}{\textbf{Training Set:}} What remains is what you use to train the models (N.B. \alert{It must be representative of what you want to learn!})
    \end{itemize}

    \vfill \pause

    {\footnotesize
        \begin{tabular}{@{}llp{28em}@{}}
            \textbf{Glossary:} &                         &                                                                                                                                     \\
                               & \emph{testing:}         & It measures the \textbf{generalisation} capability of a ML model. \alert{Do not touch this data until you give your final results!} \\
                               & \emph{validation:}      & Scores the performance of the models on a distribution similar to that used in training                                             \\
                               & \emph{model selection:} & Choice of the \textbf{model} and \textbf{hyperparameters} that deliver the best results on validation                               \\
        \end{tabular}
    }
\end{frame}

\begin{frame}{Preparing a ML model}
    Let $\mathcal{D}$ be the whole data set.
    The procedure we just saw leaves you with:
    \begin{equation*}
        \mathcal{D} = \mathcal{D}_{\text{train}} \sqcup \mathcal{D}_{\text{val}} \sqcup \mathcal{D}_{\text{test}}
    \end{equation*}
    where
    \begin{enumerate}
        \item $x \in \mathcal{D}_{\text{test}}$ represent the real \textbf{use case} scenario
        \item $x \in \mathcal{D}_{\text{val}}$ come from the same distribution as $\mathcal{D}_{\text{train}}$ and are used for evaluation
        \item $x \in \mathcal{D}_{\text{train}}$ are \textbf{shuffled} data representing \textbf{realistic}
    \end{enumerate}
    This is known as \textbf{holdout validation}.

    \vfill

    {\footnotesize
        Once a \textbf{test set} is selected, there are other types of \textbf{validation} procedures worth knowing:
        \begin{itemize}
            \item \textbf{cross-validation}: Instead of choosing one validation set, you iteratively retain a certain number of data from the train distribution for validation and train on the others
            \item \textbf{leave-one-out} (\emph{jackknife}): Train on all but one piece of data and evaluate on the remaining point. Repeat for all samples at your disposal (extreme case of cross-validation).
        \end{itemize}
    }

    \vfill

    {\tiny Clearly, the more $\mathcal{D}_{\text{train}}$ \enquote{looks like} $\mathcal{D}_{\text{test}}$ the better!}
\end{frame}
